\section{Формальная постановка задачи \textcolor{red}{и предположения}}
\label{sec:framework}

\textcolor{blue}{
В этой главе формализуются основные понятия и обозначения, используемые
в дальнейшем. Сначала
уточняется постановка задачи и вводятся обозначения,
затем для формального определения долгосрочного эффекта 
вводится понятие ограниченной пожизненной ценности (limited LTV, lLTV),
после чего показывается как эти конструкции 
применяются для А/Б эксперимента.}
% Комментарий: формулировка корректна, но немного громоздкая.
% Можно при желании подсушить ("далее вводится понятие ..." и "затем показывается применение..."),
% но это не ошибка, скорее стилистика.

Ключевая идея подхода состоит в том, что долгосрочное воздействие
экспериментального изменения проявляется не столько в смене поведения
\textit{внутри} орбиталей, сколько в \textit{перераспределении пользователей
между орбиталями}. Орбитали задают устойчивые \textcolor{red}{архетипы поведения}, а сам
эффект эксперимента реализуется за счёт того, что доли пользователей
\textcolor{red}{принаджелащих конкретным архетипам} меняются. Это допущение является центральным для
предлагаемой схемы и \textcolor{red}{в экспериментальной части работы} будет
поддержано эмпирическими наблюдениями.
% Комментарий: "в экспериментальной части работы" — нормальная ссылка,
% но для более строгого стиля можно уточнить ("в главах, посвящённых экспериментальному исследованию").
% Содержательно всё ок.

Второе важное отличие нашего подхода состоит в \textcolor{red}{выборе объекта} оценивания.
% Комментарий: термин "выбор объекта оценивания" немного расплывчат.
% Альтернативы: "в постановке задачи оценивания" или "в определении целевого объекта оценивания".
% Но с точки зрения смысла ошибок нет.
Мы стремимся оценить не гипотетический эффект бесконечно продолжительного
воздействия, а тот долгосрочный эффект, который \textit{уже сформировался}
к моменту окончания эксперимента. Такая цель согласуется с типичными
ограничениями на длительность А/Б‑экспериментов и проще реализуема
на практике, хотя и не устраняет всех проблем краткосрочного окна
(например, эффекта новизны).

\subsection{Орбитали как целеориентированные сегменты}
\label{subsec:orbitals_def}

В \textcolor{red}{нашем подходе} поведение пользователей описывается через
\textit{орбитали} — поведенческие сегменты, построенные с учётом
целевой метрики \textcolor{red}{(LTV или её прокси)} и определённые
интерпретируемыми правилами над признаками пользователя. Сегментация
является \textit{целеориентированной} (target‑aware): орбитали
конструируются так, чтобы хорошо предсказывать будущую ценность,
а не только описывать текущие различия пользователей.
% Комментарий: тезис "а не только описывать текущие различия" — риторически ок,
% но можно уточнить, что "в отличие от чисто описательных сегментаций" — если захочешь избежать двусмысленности.

Формально множество орбиталей обозначим через $\mathcal{C}$. Для
каждого пользователя $u$ и момента времени $t$ определена его
орбитальная принадлежность $c_u(t) \in \mathcal{C}$. Правила,
определяющие орбитали, считаются зафиксированными во времени и
зависят только от признаков пользователя; это позволяет вычислять
$c_u(t)$ в разные моменты, не изменяя саму схему сегментации.
% Комментарий: предположение "зависят только от признаков пользователя" — сильное.
% Формально это и есть допущение о стабильности сегментации, возможно, стоит явно назвать его предположением.

\textcolor{red}{
Конкретный алгоритм построения орбиталей (обучение модели,
извлечение правил с помощью Monoforest, семантическая разметка
сегментов) подробно описан в отдельной главе, посвящённой методу
Orbitals. В рамках настоящей главы мы будем считать, что такая
сегментация уже выполнена и множество орбиталей $\mathcal{C}$
задано.}

\subsection{Limited LTV и долгосрочный эффект}
\label{subsec:lltv}

Цель онлайн-экспериментов часто формулируется как оценка
\textit{долгосрочного эффекта} воздействия на ценность пользователя.
Прямое измерение пожизненной ценности (LTV) затруднено из-за длинных
горизонтов наблюдения и высокой дисперсии, поэтому используется понятие
\textcolor{red}{\textit{ограниченной}} пожизненной ценности (limited LTV, lLTV), задаваемой
на фиксированном прогнозном интервале.
% Комментарий: формулировка корректна; "часто формулируется" — мягко, не категорично.

Для пользователя $u$ и момента времени $t$ обозначим через $r_u(t)$
выручку, сгенерированную пользователем в момент $t$. Пусть
\[
  \mathcal{F}(t, T_F) = (t, t + T_F]
\]
— прогнозный горизонт длины $T_F$, начинающийся в момент $t$.
Тогда lLTV пользователя $u$ в момент $t$ определяется как
\begin{equation}
    R_u(t) = \sum_{t' \in \mathcal{F}(t, T_F)} r_u(t').
\end{equation}

\textcolor{red}{
Для множества пользователей $P$ сумма lLTV равна
\begin{equation}
    R(t; P) = \sum_{u \in P} R_u(t).
\end{equation}
}

% \subsection{Expected limited LTV через орбитали}

На практике точное предсказание абсолютных значений $R_u(t)$
затруднено из-за сезонности, конкурентной среды и других внешних
факторов. Во многих прикладных задачах достаточно оценить не абсолютную,
а относительную прибавку в ожидаемой выручке, которую обеспечивает
рассматриваемое изменение.
% Комментарий: утверждение про "достаточно" — скорее мотивация, чем строгий факт,
% но на уровне введения это нормально.

Ожидаение lLTV случайного пользователя из множества $P$
обозначим через
\begin{equation}
    V(t; P) = \mathbb{E}\big[ R_u(t) \,\big|\, u \sim P \big].
\end{equation}

Чтобы связать lLTV с орбиталями, рассмотрим распределение
пользователей множества $P$ по орбиталям в момент времени $t$.
Обозначим через $D(t; P)$ соответствующее распределение, а через
\begin{equation}
    p(c \mid D(t; P)) =
      \mathbb{P}\big( c_u(t) = c \,\big|\, u \sim P \big)
\end{equation}
— вероятность того, что случайно выбранный пользователь из $P$
\textcolor{red}{принадлежит} орбитали $c$.

Для каждой орбитали $c \in \mathcal{C}$ определим условное ожидание lLTV:
\begin{equation}
    V_c(t) = \mathbb{E}\big[ R_u(t) \,\big|\, c_u(t) = c \big].
\end{equation}
Тогда ожидание lLTV по множеству $P$ может быть записано как
\begin{equation}
    V(t; P) =
      \sum_{c \in \mathcal{C}} V_c(t)\, p(c \mid D(t; P)).
\end{equation}

Ключевое структурное предположение подхода состоит в том, что
относительный потенциал орбиталей по выручке стабилен во времени.
Формально предполагается, что для любых $t$ и любой пары орбиталей
$c, c'$ выполнено
\begin{equation}
    \log V_c(t) - \log V_{c'}(t) = \kappa_{c,c'},
\end{equation}
где $\kappa_{c,c'}$ — константа, зависящая только от пары орбиталей.
Это эквивалентно тому, что существуют положительные множители $V_c > 0$ и общий
временной множитель $V_o(t) > 0$ такие, что
\begin{equation}
    V_c(t) = V_c \cdot V_o(t)
    \quad \text{для всех } c \in \mathcal{C},\, t.
\end{equation}
% Комментарий: это сильное предположение. Всё правильно выписано,
% но в тексте можно (если захочешь) явно проговорить, что далее оно служит модельным допущением,
% а не проверяемым фактом.

В этой мультипликативной схеме $V_o(t)$ отражает общие временные
факторы, одинаково влияющие на все орбитали, а $V_c$ — устойчивые
различия в lLTV между орбиталями. Подобная декомпозиция
согласуется со стандартными мультипликативными моделями временных
рядов~\cite{hyndman2018forecasting} и в рассматриваемой задаче
поддерживается эмпирическими наблюдениями \textcolor{red}{(пикрил)}.

При таких обозначениях
\begin{equation}
  V(t; P)
  = \sum_{c \in \mathcal{C}} V_c(t) \, p(c \mid D(t; P))
  = V_o(t) \sum_{c \in \mathcal{C}} V_c \, p(c \mid D(t; P)).
\end{equation}

Рассмотрим два распределения $D_1$ и $D_2$. Относительная разница
в ожидании lLTV между ними равна
\begin{equation}
\label{eq:rel_diff}
r_V(D_2 \parallel D_1)
= \frac{V(t; D_2) - V(t; D_1)}{V(t; D_1)}
= \frac{\sum_c V_c \, p(c \mid D_2) - \sum_c V_c \, p(c \mid D_1)}%
       {\sum_c V_c \, p(c \mid D_1)}.
\end{equation}
Общий временной множитель $V_o(t)$ сокращается, поэтому относительный
эффект определяется только различием в составе орбиталей и их ценности
$V_c$. Это соотношение лежит в основе оценки долгосрочного эффекта
через перераспределение пользователей между орбиталями.

\subsection{Переходы между орбиталями в А/Б-экспериментах}
\label{subsec:ab_orbitals}

В постановке А/Б-эксперимента требуется сравнить  lLTV
тестовой и контрольной групп.
% Комментарий: здесь "требуется сравнить lLTV" чуть телеграфно.
% Можно "требуется сравнить ожидаемую lLTV тестовой и контрольной групп" — для симметрии с \(V(t;P)\).

Будем считать, что обе \textcolor{red}{популяции}
достаточно велики и в теории \textcolor{red}{являются несмещёнными}, так что в момент начала
эксперимента $T_S$ распределения пользователей по орбиталям 
\textcolor{red}{$D_A(T_S)$ и $D_B(T_S)$ совпадают}.


\textcolor{red}{Пусть $T_E$ — момент окончания эксперимента. Тогда интервал $[T_S, T_E]$ 
задаёт период действия экспериментального воздействия, а lLTV оценивается на прогнозном горизонте}
\[
  \mathcal{F}(T_E, T_F) = (T_E, T_E + T_F].
\]
% Комментарий: фраза "задаёт период действия экспериментального воздействия" корректна,
% но можно чуть суше ("будем называть окном эксперимента").

Согласно формуле~\eqref{eq:rel_diff}, для сравнения групп можно
использовать относительное различие в ожидании lLTV. Однако вместо
непосредственного сравнения орбитальных распределений в момент $T_E$
удобно рассматривать их изменение за время эксперимента. Для этого
вводится скорректированная метрика
\begin{equation}
\label{eq:rel_diff_corrected}
\begin{split}
r_{V\text{-corr}}(D_B \parallel D_A)
=
\frac{\big(V(T_E; D_B) - V(T_S; D_B)\big)
-
\big(V(T_E; D_A) - V(T_S; D_A)\big)}
{V(T_E; D_A)}.
\end{split}
\end{equation}

Такое преобразование позволяет использовать изменение распределений
за время эксперимента, а не только их значения в конечной точке. Этот
приём уменьшает шум, связанный со стартовой точкой эксперимента. Хотя
в теории при корректной рандомизации распределения $D_A(T_S)$ и
$D_B(T_S)$ совпадают, на практике орбитальное распределение в начальный
момент является случайным вектором, и потому между контрольной и
тестовой группами могут возникать небольшие различия. Скорректированная
метрика частично компенсирует эти различия.
% Комментарий: логика согласована с оригинальной статье; формулировки ок.

С точки зрения поведения пользователей это означает анализ переходов
между орбиталями \textcolor{red}{на интервале $[T_S, T_E]$}. Для каждой группы
$g \in \{A,B\}$ можно рассматривать эмпирическую матрицу переходов
\[
  P^{(g)}_{c \to c'} =
  \mathbb{P}\big( c_u(T_E) = c' \,\big|\, c_u(T_S) = c,\ u \in P_g \big),
\]
которая описывает вероятность перехода из орбитали $c$ в орбиталь $c'$
за время эксперимента. Различия этих матриц между группами приводят
к различиям в итоговых распределениях $D_A(T_E)$ и $D_B(T_E)$, а значит,
и в значениях $r_V$ и $r_{V\text{-corr}}$.

Таким образом, в рамках \textcolor{red}{нашего подхода} задача оценки долгосрочного
эффекта в А/Б-эксперименте сводится к сравнению того, как экспериментальное
воздействие изменяет распределение пользователей по орбиталям и как эти
изменения затем отражаются на ожидании lLTV.
% Комментарий: это резюме уже звучало по смыслу выше (во введении),
% но дублирование в конце подглавы вполне допустимо.

\subsection{\textcolor{red}{Поставленные задачи оценки и эталонные методы}}

В рамках предложенного подхода оценки естественным образом
возникают две основные задачи:

\begin{enumerate}
    \item \textbf{Присвоение орбиталей.}
    Необходимо построить сегментацию пользователей на орбитали по их
    историческим признакам так, чтобы принадлежность к орбитали
    отражала устойчивые поведенческие паттерны, предсказывающие
    будущую выручку. Иными словами, орбиталь должна служить
    компактным суррогатом индивидуального поведения пользователя
    с точки зрения lLTV.

    \item \textbf{Оценка ценности орбиталей.}
    Требуется получить оценки инвариантных по времени множителей $\hat{V}_c$, 
    характеризующих относительную ценность каждой орбитали $c \in \mathcal{C}$ с точки
    зрения ожидаемой lLTV. Эти множители связывают
    поведенческие сегменты с их вкладом в будущую выручку.
\end{enumerate}

При наличии орбитальной принадлежности пользователей и набора
$\{\hat{V}_c\}_{c \in \mathcal{C}}$ можно вычислять метрики $r_V$ и
$r_{V\text{-corr}}$ по наблюдаемым распределениям пользователей по
орбиталям в начале и в конце эксперимента. Таким образом, относительный
эффект по lLTV оказывается выражен через изменения
распределения по орбиталям и их ценности, что позволяет оценивать
долгосрочное воздействие эксперимента, не дожидаясь полного завершения
прогнозного горизонта $\mathcal{F}(T_E, T_F)$.

Для оценки качества орбитального оценивателя рассмотрим два
естественных эталонных подхода.


\paragraph{(1) Прямое наблюдение ограниченной LTV.}

Наиболее прямой способ оценки долгосрочного эффекта состоит в том,
чтобы дождаться, пока прогнозное окно $\mathcal{F}(T_E, T_F)$ полностью
завершится, и затем непосредственно вычислить реализованную ограниченную
LTV каждого пользователя:
\[
  R_u(T_E) = \sum_{t' \in \mathcal{F}(T_E, T_F)} r_u(t').
\]
Далее средние значения $R_u(T_E)$ в контрольной и тестовой группах
используются для оценки относительного эффекта.


В идеализированной постановке такой оцениватель является несмещённым:
он опирается на фактические денежные потоки и не зависит от
модельных предположений. Однако на практике это свойство может
нарушаться, если за длительный период наблюдения поведение
пользователей в группах меняется по‑разному (например, из‑за
перекрёстных запусков других экспериментов, маркетинговых активностей
или сдвига состава аудитории). Кроме того, прямое наблюдение
ограниченной LTV страдает от высокой дисперсии, поскольку распределение
выручки по пользователям обычно имеет тяжёлые хвосты.


\paragraph{(2) Сильная предсказательная модель.}


Альтернативный эталонный подход состоит в построении
супервизируемой модели, предсказывающей $R_u(T_E)$ по признакам
пользователя, наблюдаемым к моменту $T_E$. Такая модель стремится
максимально использовать информацию о текущем состоянии пользователя
(частота и давность активности, структура трат, признаки вовлечённости
и т.п.), чтобы оценить его будущую ограниченную LTV.

Если модель хорошо специфицирована и обучена на репрезентативных
данных, она даёт оцениватель с существенно меньшей дисперсией по
сравнению с прямым наблюдением: индивидуальные шумовые колебания
денежных потоков сглаживаются за счёт агрегации информации в признаках.
Однако на практике такой подход чувствителен к ряду факторов:
ошибкам спецификации модели, сдвигу распределений признаков между
обучением и применением (covariate shift), переобучению и другим
источникам систематического смещения. В результате оценки эффекта
лечения могут быть искажены даже при визуально хорошем качестве
предсказаний.


\paragraph{Орбитальная оценка.}

Предлагаемый орбитальный оцениватель на основе метрик $r_V$ и
$r_{V\text{-corr}}$ занимает промежуточное положение между этими двумя
крайними подходами. С одной стороны, он опирается на структурное
модельное предположение
\[
  V_c(t) = V_c \cdot V_o(t),
\]
которое в реальных данных может выполняться лишь приближённо. Если
это предположение нарушается, оценки, основанные на орбиталях, могут
наследовать небольшое смещение.

С другой стороны, поведение пользователей сводится к относительно
небольшому числу орбиталей, и требуется оценить лишь низкоразмерный
вектор параметров $\{V_c\}_{c \in \mathcal{C}}$. Это радикально
снижает размерность задачи по сравнению с полностью индивидуальными
предсказаниями $R_u(T_E)$ и позволяет существенно уменьшить дисперсию
оценок эффекта. По сути, орбитали играют роль агрегированных суррогатов:
они достаточно грубо аппроксимируют сложное пространство
поведенческих траекторий, но при этом сохраняют различия между
типичными режимами поведения, важные с точки зрения будущей ценности.

Таким образом, орбитальный оцениватель реализует осознанный компромисс
между небольшим модельным смещением и существенным снижением дисперсии.
В контексте A/B‑экспериментов, где статистическая мощность часто
ограничивает возможность обнаружения слабых, но практически значимых
эффектов, такой компромисс оказывается особенно важен: снижение
дисперсии метрики непосредственно повышает чувствительность теста
при фиксированном объёме выборки.\textcolor{red}{[file:97]}


В \textcolor{red}{экспериментальной части работы} этот компромисс будет
квантифицирован на реальных данных: будут сопоставлены прямое
наблюдение ограниченной LTV, сильная предсказательная модель и
орбитальный оцениватель по смещению, дисперсии и способности
корректно выявлять эффект в A/A‑ и A/B‑экспериментах.